\documentclass[a4paper,11pt]{article}
\usepackage{exam-style-842-unit}
\examinehead{842 信息化概论 · 第 4 单元练习册}{练习二 参考答案与解析}
\begin{document}

\unitanswerscover{842 信息化概论 · 单元练习册}{第 4 单元参考答案}{练习二（开发架构与 Web 服务）的答案与解析；每道判断题均给出改正后的正确表述。说明：教材未设 B/S、C/S 与 Web 服务专节，本册按考试大纲“二、4”与通用规范知识补充，相关题目标注“教材未覆盖，按大纲补充”。}

\noindent{\bfseries\color{ctp-mauve}【关于教材依据的说明】}\textbf{教材未覆盖，按大纲补充。}经对《信息技术导论》OCR 全文检索确认：教材第 3 章 3.3 节只讲软件工程、软件过程（瀑布模型）与软件项目管理，\textbf{没有} B/S、C/S 架构与 Web 服务的专节；全书仅有三处零星线索——3.2.4 讲 DB2 时提到“客户机/服务器用户”（第 36 页）、6.3.1 讲 TCP/IP 应用层时列出 HTTP（第 108 页）、4.1.1 讲云计算由来时提到亚马逊提供“Web Service”（第 45 页）。因此练习二的架构与 Web 服务部分按 2027 年考试大纲“二、计算机与软件技术”第 3、4 条（软件工程基本概念、软件开发主要流程、软件项目管理；B/S、C/S 等软件开发主要架构，Web 服务及其体系架构）与计算机领域通用规范知识命题，属考纲要求而教材未展开的内容。\par\vspace{0.6em}

\answerpart{一、单项选择题（10 题，每题 3 分，共 30 分）}

\q B。\begin{analysis}B/S（Browser/Server，浏览器/服务器）架构以 Web 浏览器作为统一客户端，业务逻辑与数据集中在服务器端，客户端零安装、升级维护集中在服务器，适合互联网环境下的大量用户访问（教材未覆盖，按大纲补充）。A C/S 架构需在客户端安装专用程序；C 单机架构不涉及网络分工；D 主机—终端架构由主机集中处理、终端只做输入输出，与 B/S 的“浏览器即客户端”不同。\end{analysis}

\q B。\begin{analysis}C/S（Client/Server，客户机/服务器）架构中，客户端安装专用程序，承担表示逻辑与部分业务逻辑，服务器（通常是数据库服务器）承担数据管理与共享资源服务，因此是“胖客户端”（教材未覆盖，按大纲补充）。A 只负责显示的是 B/S 中的浏览器角色；C 数据管理由服务器承担；D 网络路由由网络设备完成。\end{analysis}

\q A。\begin{analysis}B/S 架构最突出的优点是客户端零安装、零维护，系统升级只需更新服务器端，用户用浏览器即可使用，且跨平台、易扩展（教材未覆盖，按大纲补充）。B 与 B/S 的优点相反，是 C/S 的特征；C 说法错误，B/S 同样受网络带宽和延迟影响；D 错误，业务逻辑主要运行在服务器端，浏览器只负责表示与交互。\end{analysis}

\q B。\begin{analysis}三层架构中：表示层负责与用户交互、信息展示与输入校验；业务逻辑层负责业务规则与业务流程处理，是系统的核心；数据访问层负责封装数据库访问操作。故 B 正确（教材未覆盖，按大纲补充）。A 是表示层职责；C 是数据访问层职责；D 是网络层的职责。\end{analysis}

\q B。\begin{analysis}Web 服务三要素中，WSDL（Web Services Description Language）用 XML 描述服务的接口、操作、消息格式与访问地址，相当于“服务说明书”；SOAP 负责消息传递；UDDI 负责服务的注册与发现；HTTP 是传输协议（教材未覆盖，按大纲补充）。\end{analysis}

\q C。\begin{analysis}UDDI（Universal Description, Discovery and Integration）提供服务注册、发布与查找机制，使服务请求者能够发现所需服务，相当于“服务黄页”（教材未覆盖，按大纲补充）。A SOAP 是消息协议；B WSDL 是服务描述；D SSL 是安全传输协议，均不负责注册与发现。\end{analysis}

\q B。\begin{analysis}SOAP（Simple Object Access Protocol，简单对象访问协议）采用 XML 编码消息，通常绑定 HTTP 传输，从而实现跨平台、跨语言的服务调用（教材未覆盖，按大纲补充）。A 二进制、C JSON 是其他接口形式（如 RPC、REST/JSON）常用的表示，不是 SOAP 的标准编码；D HTML 是网页标记语言，不是消息编码格式。\end{analysis}

\q B。\begin{analysis}在浏览器地址栏输入网址后，浏览器先解析 URL（协议、主机、端口、路径等），再通过 DNS 把域名解析为 IP 地址，然后与服务器建立 TCP 连接并发送 HTTP 请求，最后才解析渲染页面（教材未覆盖，按大纲补充；教材 6.3.1 第 108 页在 TCP/IP 应用层列出 HTTP，可作协议归属依据）。A 未完成域名解析无法确定服务器地址；C 渲染发生在收到响应之后；D 建立数据库连接是服务器端应用的行为，不是浏览器最先进行的操作。\end{analysis}

\q B。\begin{analysis}HTTP 状态码用于表示服务器对请求的处理结果：200 表示请求成功，301/302 表示重定向，400 表示请求错误，403 表示禁止访问，404 表示请求的资源不存在，500 表示服务器内部错误（教材未覆盖，按大纲补充；可结合教材 6.3.1 应用层协议理解 HTTP 的请求—响应模式，第 108 页）。\end{analysis}

\q C。\begin{analysis}Web 服务的特点是跨平台、跨语言、基于标准协议、松耦合、以服务为单位封装业务功能，可在互联网范围内调用（教材未覆盖，按大纲补充）。C“只能在同一个局域网内调用”与 Web 服务“跨网络、跨组织集成”的初衷相反，故 C 符合题意。\end{analysis}

\answerpart{二、多项选择题（3 题，每题 4 分，共 12 分）}

\q ABCD。\begin{analysis}B/S 架构相对 C/S 的优点包括：客户端零安装、零维护；升级维护集中在服务器端；跨平台、易扩展；适合在互联网上向大量用户提供服务（教材未覆盖，按大纲补充）。A、B、C、D 四项均正确。注意 B/S 的不足是交互性与响应速度受网络影响、对浏览器兼容性有依赖、复杂图形与大数据量实时处理能力相对较弱。\end{analysis}

\q ABC。\begin{analysis}Web 服务的三要素是 SOAP（消息协议）、WSDL（服务描述）、UDDI（注册与发现）；此外 Web 服务还依赖 HTTP、XML 等标准协议与数据格式（教材未覆盖，按大纲补充）。D HTML 是网页标记语言，用于浏览器展示，不属于 Web 服务三要素。\end{analysis}

\q ABCD。\begin{analysis}网页请求与响应过程包括：浏览器解析 URL 并进行 DNS 域名解析；与 Web 服务器建立 TCP 连接（三次握手）；发送 HTTP 请求报文（请求行、请求头、请求体）；服务器处理请求（静态资源直接返回，动态请求交由应用服务器并访问数据库）；返回 HTTP 响应报文（状态行、响应头、响应体）；浏览器解析 HTML 并请求 CSS、JS、图片等资源后渲染页面（教材未覆盖，按大纲补充）。四项均可能发生。\end{analysis}

\answerpart{三、判断题（6 题，每题 2 分，共 12 分）}

\q √。\begin{analysis}B/S 架构以浏览器为统一客户端，用户无需安装专用客户端软件，系统升级只需更新服务器端，因而部署快、维护成本低、易扩展（教材未覆盖，按大纲补充）。\end{analysis}

\q ×。改正：C/S 架构的客户端安装了专用程序并承担部分业务逻辑，系统升级时客户端程序也需要安装或更新，维护成本较高，不能只升级服务器。\begin{analysis}C/S 的“胖客户端”特征决定了业务逻辑与界面在客户端有分布，客户端与服务器需要配套升级；这正是 B/S 架构要解决的问题之一（教材未覆盖，按大纲补充）。\end{analysis}

\q ×。改正：Web 服务基于标准 Web 协议与中立的 XML 数据格式，可以实现跨平台、跨语言的应用之间互相调用。\begin{analysis}Web 服务的核心价值就是屏蔽操作系统、编程语言与硬件平台的差异，以服务方式实现松耦合集成，因此“只能在同一语言、同一平台内调用”的说法与 Web 服务的定义相反（教材未覆盖，按大纲补充）。\end{analysis}

\q √。\begin{analysis}SOAP 是基于 XML 的消息协议，规定了消息的结构与处理规则，通常绑定 HTTP 传输，从而穿越防火墙、便于在互联网上调用（教材未覆盖，按大纲补充）。\end{analysis}

\q ×。改正：WSDL 负责描述服务的接口、操作与访问地址；服务的注册与发现由 UDDI 完成。\begin{analysis}三要素分工不同：SOAP 管“怎么传消息”，WSDL 管“服务长什么样、怎么调”，UDDI 管“到哪里找服务”（教材未覆盖，按大纲补充）。\end{analysis}

\q √。\begin{analysis}HTTP 响应的响应体可以是 HTML 页面，也可以是 JSON、XML、图片、视频、文件流等数据；前端页面通过 Ajax/Fetch 请求接口时，服务器通常返回 JSON 数据而不是完整 HTML（教材未覆盖，按大纲补充）。\end{analysis}

\answerpart{四、填空题（6 空，每空 1 分，共 6 分）}

\q 浏览器/服务器（Browser/Server）。\begin{analysis}C/S 为客户机/服务器（Client/Server），B/S 为浏览器/服务器（Browser/Server），二者区别在客户端形态与业务逻辑分布（教材未覆盖，按大纲补充）。\end{analysis}

\q WSDL。\begin{analysis}Web 服务三要素：SOAP（消息协议）、WSDL（服务描述）、UDDI（注册与发现）（教材未覆盖，按大纲补充）。\end{analysis}

\q 业务逻辑层。\begin{analysis}三层架构＝表示层＋业务逻辑层＋数据访问层；业务逻辑层实现业务规则与流程，是系统的核心层（教材未覆盖，按大纲补充）。\end{analysis}

\q 请求头。\begin{analysis}HTTP 请求报文由请求行（方法、URL、协议版本）、请求头（Host、User-Agent、Content-Type 等）和请求体（实体）组成（教材未覆盖，按大纲补充）。\end{analysis}

\q 请求的资源不存在（Not Found）。\begin{analysis}404 表示服务器找不到请求的资源，常见于 URL 拼写错误或资源已被删除；200 表示成功，500 表示服务器内部错误（教材未覆盖，按大纲补充）。\end{analysis}

\q DNS（域名）。\begin{analysis}DNS（Domain Name System，域名系统）把便于记忆的域名解析为 IP 地址；教材 6.3.1（第 108 页）在 TCP/IP 应用层协议中列出了 DNS，可作对应依据。\end{analysis}

\answerpart{五、简答题（4 题，每题 6 分，共 24 分）}

\q 参考答案要点（教材未覆盖，按大纲补充）：
（1）C/S 架构（1.5 分）：即客户机/服务器架构，客户端安装专用程序，承担表示逻辑与部分业务逻辑，服务器承担数据管理与共享资源服务，典型为两层结构（客户端＋数据库服务器）。
（2）B/S 架构（1.5 分）：即浏览器/服务器架构，以 Web 浏览器作为统一客户端，业务逻辑与数据集中在服务器端，通常采用表示层、业务逻辑层、数据访问层三层结构。
（3）主要区别（2 分）：客户端形态——C/S 需安装专用客户端，B/S 只需浏览器；业务逻辑分布——C/S 客户端承担部分逻辑，B/S 集中在服务器；升级维护——C/S 需客户端与服务器配套升级，B/S 只升级服务器；跨平台性——B/S 更强；交互与性能——C/S 交互性强、响应快，B/S 受网络影响较大。
（4）适用场景（1 分）：C/S 适合局域网内、交互复杂、数据量大、实时性要求高的业务系统（如大型进销存、工业控制）；B/S 适合互联网上面向大量用户的查询、发布、电商与在线办事类应用。

\q 参考答案要点（教材未覆盖，按大纲补充）：
（1）定义（2 分）：Web 服务是一种面向服务的分布式计算技术，它基于 HTTP 等标准 Web 协议，以 XML 等中立格式交换消息，使不同平台、不同语言的应用能够以松耦合方式互相调用，是 SOA 的主要实现技术。
（2）三种角色（2 分）：服务提供者（发布并实现服务）、服务注册中心（登记与查找服务，对应 UDDI）、服务请求者（查找服务并绑定调用）。
（3）三项核心技术（2 分）：SOAP——基于 XML 的消息协议，负责服务之间的消息传递；WSDL——用 XML 描述服务接口、操作、消息与访问地址；UDDI——提供服务注册、发布与发现机制。三者的关系可概括为：用 WSDL 描述服务，用 UDDI 发布和查找服务，用 SOAP 传递调用消息。

\q 参考答案要点（教材未覆盖，按大纲补充）：
（1）表示层（1.5 分）：直接面向用户，负责界面展示、输入采集与基本校验，常见形式为 Web 页面、移动 App、小程序等。
（2）业务逻辑层（1.5 分）：实现业务规则、业务流程与业务服务，是系统的核心，接收表示层请求、调用数据访问层。
（3）数据访问层（1.5 分）：封装对数据库、文件等数据源的访问，为上层提供统一的数据操作接口。
（4）分层的主要好处（1.5 分）：职责分离、层间松耦合，便于分工开发与维护；某一层变化时对其它层影响小，可复用、可扩展、可替换；便于测试与安全管理。三层架构是 B/S 应用的典型结构，C/S 也可采用三层结构以减轻客户端负担。

\q 参考答案要点（教材未覆盖，按大纲补充；HTTP 归属可参见教材 6.3.1，第 108 页）：
（1）解析 URL（1 分）：浏览器解析用户输入的 URL，得到协议（http/https）、主机域名、端口、路径与查询参数。
（2）DNS 域名解析（1 分）：通过 DNS 把域名解析为服务器的 IP 地址。
（3）建立连接（1 分）：浏览器与 Web 服务器建立 TCP 连接（三次握手）；若为 HTTPS，还需进行 TLS 握手。
（4）发送 HTTP 请求（1 分）：浏览器发送请求报文，包含请求行（方法如 GET/POST、URL、协议版本）、请求头和请求体。
（5）服务器处理并响应（1 分）：Web 服务器接收请求，静态资源直接返回；动态请求转交应用服务器，由业务逻辑访问数据库后生成结果，返回 HTTP 响应报文（状态行、响应头、响应体）。
（6）浏览器渲染（1 分）：浏览器解析 HTML，按需请求 CSS、JavaScript、图片等资源，完成布局与绘制，呈现页面；同时执行脚本实现交互。

\answerpart{六、综合题（1 题，共 16 分）}

\q 参考答案要点（教材未覆盖，按大纲补充）：
（1）改造理由与不足（5 分）
理由（答出 3 点得 3 分）：第一，使用范围扩大——B/S 客户端零安装，合作社、门店、第三方平台人员用浏览器即可访问，不必逐台安装客户端；第二，升级维护方便——业务逻辑集中在服务器端，系统升级只需更新服务器，避免 C/S 客户端逐台升级；第三，跨平台与易扩展——浏览器天然跨操作系统，便于接入微信小程序、移动端和第三方电商平台；第四，适合互联网开放访问——B/S 基于 HTTP，便于与外部系统集成和对外提供服务。
不足（2 分，答出 2 点）：交互性与响应速度受网络带宽和延迟影响，弱网体验不如 C/S；复杂图形处理、大数据量实时计算能力相对较弱；对浏览器兼容性与前端安全性依赖较强。
（2）三层架构划分与职责（4 分，每层 1 分、总体分层合理 1 分）：表示层——面向办公人员、合作社、小程序与电商平台的多端界面，负责展示与输入校验；业务逻辑层——实现库存查询、订单提交、出入库、结算等业务规则与流程，并以 Web 服务接口对外暴露；数据访问层——封装库存、订单、商品、客户等数据表的访问，保证数据一致性与安全性。三层之间通过明确定义的接口调用，前端不直接访问数据库。
（3）三项核心技术及作用（4 分）：SOAP——定义并承载服务调用的消息格式，基于 XML、通常绑定 HTTP，实现跨平台消息交换；WSDL——用 XML 描述对外服务的接口、操作、参数与访问地址，供小程序与第三方平台开发者查阅对接；UDDI——提供服务的注册、发布与查找机制，使调用方能够发现企业的服务（企业内部也可用统一的服务目录/注册中心替代）。三者配合实现“描述—发布—调用”的完整闭环。
（4）请求与响应过程及注意事项（3 分）
过程（2 分）：小程序发起调用时，先由 URL 定位服务端点，经 DNS 解析得到服务器 IP，与服务器建立 HTTPS 连接；按 WSDL 约定构造 SOAP（或 REST/JSON）请求报文，通过 HTTP POST 发送到 Web 服务器；服务器交由应用服务器执行业务逻辑、访问数据库，处理完成后返回含状态码与响应体的 HTTP 响应；小程序解析结果并展示。
注意事项（1 分，答出 2 点即可）：安全方面——使用 HTTPS/SSL 加密传输，接口做身份认证与访问授权（如数字证书、令牌），敏感数据脱敏，防止重放与越权；性能方面——对高频库存查询做缓存、对接口做限流与并发控制，合理设计批量接口，避免小程序端轮询造成压力；同时做好接口版本管理与日志审计。

\end{document}
